\section{Przykładowe predykcje}
\label{sec:predykcje}

Aby pokazać działanie sieci na konkretnych artykułach, dla każdego tekstu ze zbioru
testowego obliczono prawdopodobieństwo $P(\text{Real})$ i porównano z progiem
0{,}5. Tabela~\ref{tab:pred} zestawia dwa pewne trafienia każdej klasy oraz oba typy
błędów: \emph{false positive} (Fake uznane za Real) i \emph{false negative} (Real
uznane za Fake).

\begin{table}[H]
  \centering
  \caption{Przykładowe predykcje MLP na zbiorze testowym (próg 0{,}5).}
  \label{tab:pred}
  \small
  \begin{tabular}{p{6.6cm}cccc}
    \toprule
    \textbf{Fragment tekstu} & \textbf{Prawdziwa} & \textbf{P(Real)} & \textbf{Predykcja} & \\
    \midrule
    \multicolumn{5}{l}{\emph{Pewne trafienia}}\\
    \addlinespace[2pt]
    ,,germany foreign minister said on saturday that\dots'' & Real & 1{,}0000 & Real & \checkmark \\
    ,,china called for ceasefire in myanmar rakhine\dots''  & Real & 1{,}0000 & Real & \checkmark \\
    ,,patrick henningsen the longer this soap opera\dots''  & Fake & 0{,}0000 & Fake & \checkmark \\
    ,,patrick henningsen despite repeated failures\dots''   & Fake & 0{,}0000 & Fake & \checkmark \\
    \addlinespace[3pt]
    \multicolumn{5}{l}{\emph{Błędy: false positive (Fake $\to$ Real)}}\\
    \addlinespace[2pt]
    ,,the washington post nearly third of territory\dots''  & Fake & 0{,}9998 & Real & $\times$ \\
    ,,leading from behind hasn worked out so well\dots''    & Fake & 0{,}9997 & Real & $\times$ \\
    \addlinespace[3pt]
    \multicolumn{5}{l}{\emph{Błędy: false negative (Real $\to$ Fake)}}\\
    \addlinespace[2pt]
    ,,- below are the highlights from oct. 25 excl\dots''   & Real & 0{,}0015 & Fake & $\times$ \\
    ,,the emmy awards show was peppered with politic\dots'' & Real & 0{,}0059 & Fake & $\times$ \\
    \bottomrule
  \end{tabular}
\end{table}

\paragraph{Kalibracja i omówienie błędów.}
Prawdopodobieństwa $P(\text{Real})$ są \textbf{nasycone na krańcach}
($\approx$0 dla Fake, $\approx$1 dla Real), również w przypadku pomyłek: artykuły
Fake błędnie zaklasyfikowane jako Real otrzymują $P(\text{Real})\approx0{,}9998$.
Model jest zatem \textbf{silny w rankingu} (ROC-AUC 0{,}9987), lecz \textbf{przepewny
i słabo skalibrowany}: myli się rzadko, ale gdy się myli, robi to z wysoką pewnością,
zamiast sygnalizować niepewność wartościami w pobliżu progu 0{,}5. Same błędy dotyczą
tekstów stylistycznie ,,pogranicznych'', co \emph{sugeruje} (lecz nie dowodzi)
związek z cechami z sekcji~\ref{sec:eval}:

\begin{itemize}[noitemsep]
  \item \textbf{False positive} (Fake $\to$ Real): teksty otwierające się
        frazą agencyjno-prasową (np.\ ,,the washington post\dots''), stylistycznie
        zbliżone do źródeł wiarygodnych;
  \item \textbf{False negative} (Real $\to$ Fake): teksty nietypowe dla klasy Real,
        sformatowane jako lista (,,- below are the highlights\dots'') lub o tematyce
        rozrywkowej (gala Emmy) zamiast politycznej/agencyjnej.
\end{itemize}

Błędy są więc rzadkie i koncentrują się na artykułach ,,pogranicznych'' stylistycznie,
co jest spójne z nielicznymi pomyłkami w macierzy pomyłek (Rys.~\ref{sec:eval}).
Wysoka pewność tych pomyłek pozostaje jednak sygnałem ograniczonej kalibracji modelu.
